\documentclass[11pt]{article}
\usepackage[a4paper,margin=18mm]{geometry}
\usepackage{fontspec}
\setmainfont{Latin Modern Roman}
\usepackage{amsmath,amssymb,amsthm,amscd,graphicx,color}
\usepackage{xurl}
\usepackage[unicode,hidelinks]{hyperref}
\allowdisplaybreaks
\setlength\emergencystretch{3em}

\makeatletter
\newtheorem{theorem}{Theorem}
\newtheorem*{theorem*}{Theorem}
\newtheorem{lemma}{Lemma}
\newtheorem*{lemma*}{Lemma}
\newtheorem{corollary}{Corollary}
\newtheorem*{corollary*}{Corollary}
\newtheorem{proposition}{Proposition}
\newtheorem*{proposition*}{Proposition}
\newtheorem{conjecture}{Conjecture}
\newtheorem*{conjecture*}{Conjecture}
{\theoremstyle{definition} \newtheorem{definition}{Definition}
\newtheorem*{definition*}{Definition}
\newtheorem{example}{Example}
\newtheorem*{example*}{Example}
\newtheorem{remark}{Remark}
\newtheorem*{remark*}{Remark}
\newtheorem{note}{Note}
\newtheorem*{note*}{Note}
}
\makeatother

% Commands

\newcommand{\norm}[1]{\left\Vert {#1} \right\Vert}
\newcommand{\abs}[1]{\left\vert {#1} \right\vert}
\newcommand{\set}[1]{\left\{ {#1} \right\}}
\newcommand{\parenthese}[1]{\left( {#1} \right)}
\newcommand{\dv}[2]{\frac{\mathrm d{#1}}{\mathrm d{#2}}}
\newcommand{\pd}[2]{\frac{\partial{#1}}{\partial{#2}}}
\newcommand{\setN}{{\mathbb N}}
\newcommand{\setZ}{{\mathbb Z}}
\newcommand{\setR}{{\mathbb R}}
\newcommand{\setC}{{\mathbb C}}
\newcommand{\implie}{\Rightarrow}
\newcommand{\Limplie}{\ \ \Longrightarrow\ \ }
\newcommand{\equi}{\Leftrightarrow}
\newcommand{\Lequi}{\Longleftrightarrow}
\newcommand{\precc}{\prec\!\!\prec}

\begin{document}
\begin{center}\Large On a time-oriented globally hyperbolic four-dimensional spacetime, Lorentzian distance is the infimum of positive causal-function increments under the global past-directed timelike eikonal bound\end{center}
\noindent\textbf{Stable ID:} 2784\par
\noindent\textbf{Authors:} Nicolas Franco\par
\noindent\textbf{Original work:} Global Eikonal Condition for Lorentzian Distance Function in Noncommutative Geometry\par
\noindent\textbf{Version:} Official SIGMA 6 (2010) article 064, DOI 10.3842/SIGMA.2010.064; journal PDF and native TeX acquired 2026-10-01, exact bytes pinned by SHA256\par
\noindent\textbf{Selected task scope:} Theorem1 in its fixed A=\allowbreak{}continuous causal functions setting, on an oriented and time-oriented globally hyperbolic four-dimensional spacetime with metric signature(-,+\allowbreak{},+\allowbreak{},+\allowbreak{}). For all p,q, directed Lorentzian distance d(p,q) equals inf\_f max(0,f(q)-f(p)), where admissible continuous causal f have the source almost-everywhere defined gradient, essSup\_M g(grad f,grad f)<=\allowbreak{}-1 and past-directed gradient. No replacement of A by arbitrary Sobolev/absolutely continuous or smooth test functions, no noncommutative/operator distance formula, and no removal of global hyperbolicity.\par
\noindent\textbf{Difficulty:} H2: The reverse Cauchy-Schwarz inequality gives only one direction of the distance formula; obtaining the infimum requires global functions approximating a prescribed pair of points while satisfying a timelike bound everywhere. The author builds locally finite causal-cone covers of a Cauchy hypersurface and sums Lorentzian point-distance functions outside the protected region. Separate comparable and incomparable constructions combine past and future support controls to recover both the positive distance and its zero branch.\par
\noindent\textbf{Editorial boundary note (not author text):} The complete original Theorem1 is retained in its fixed continuous causal-function class on oriented/time-oriented globally hyperbolic four-dimensional spacetimes. All new locally finite causal-cone covers, protected-region point-distance sums, past/future gradient estimates and comparable, reversed and incomparable point-pair constructions are bundled. Standard Lorentzian cut-locus/eikonal properties and Cauchy-surface existence remain cited prerequisites. The original monotone path inequality and the later caveats about other function spaces are preserved; no smooth-test, Sobolev or noncommutative extension is selected.\par
\noindent\textbf{Source:} \url{https://sigma-journal.com/2010/064/}\quad DOI: \url{https://doi.org/10.3842/SIGMA.2010.064}\par
\noindent\textbf{License and attribution:} Copyright retained by the named authors. Published by SIGMA under CC BY 4.0: \url{https://creativecommons.org/licenses/by/4.0/}. Source: \url{https://sigma-journal.com/about.html}. Adaptation consists of standalone excerpt formatting, cover and numbering restoration; original author language and mathematical text are retained.\par
\noindent\textbf{Editorial symbol notice (not author text):} The separate Notes10 extension to other function spaces requests additional absolute-continuity regularity of point-distance functions. The selected claim retains the original continuous causal-function class and its stated monotone path inequality; no extra extension or repaired regularity argument is supplied.\par
\medskip\hrule\medskip\noindent\textbf{Author text begins below.}\par
\setcounter{section}{2}
\setcounter{subsection}{0}
\setcounter{subsubsection}{0}
\setcounter{equation}{0}
% Original source lines 111--465
\section{Lorentzian distance and the timelike eikonal inequality}

From this point, we will switch from Riemannian manifolds to Lorentzian ones. We will assume that the reader is familiar with Lorentzian geometry, so we will only give def\/initions and pro\-per\-ties which are non usual or of particular interest. For the remaining, the reader can refer to~\cite{Beem,ONeill}. We will use $p \preceq q$ to indicate that $p = q$ or there is a future directed causal curve from~$p$ to~$q$, and $p \precc q$ to indicate that there is a future directed timelike curve from~$p$ to~$q$. $p \prec q$ means that $p \preceq q$ and $p\neq q$. The signature of the metric will be $(-,+,+,+)$.

\begin{definition}
The Lorentzian distance function on an oriented Lorentzian manifold $(M,g)$ is the function $d : M \times M \rightarrow [0,+\infty) \cup \set{+\infty}$ def\/ined by
\begin{gather*}
d(p,q) =
\begin{cases}
 \sup\{ l(\gamma) : \gamma \text{ future directed causal piecewise} \ & \ \\
  \qquad \text{smooth curve from $p$ to $q$}\} &\text{if } p \prec q,\\
 0 &\text{if } p \nprec q,
\end{cases}
 \end{gather*}
 with $l(\gamma) = \int \sqrt{ -g( \dot\gamma(t) ,  \dot\gamma(t) )}\ dt$ the length of the curve.
\end{definition}

Lorentzian geometry is really dif\/ferent from Riemannian one on lots of aspects. For example, the distance between to points is no more commutative, and is in fact a function which is not null only on a single future cone. Any formula for Lorentzian distance must respect these aspects. Practically, taking two arbitrary distinct points $p$ and $q$, we have to dissociate those three cases:
\begin{itemize}\itemsep=0pt
\item $p \prec q$, where $p$ and $q$ are causally related, with $q$ in the future of $p$;
\item $q \succ p$, where $p$ and $q$ are causally related, with $q$ in the past of $p$;
\item $p \nprec q$ and $q \nprec p$, where $p$ and $q$ are not causally related.
\end{itemize}

Moreover, usual inequalities as triangle or Schwartz do not hold in Lorentzian geometry, but we have inverse counterparts:

\begin{lemma}
In Lorentzian geometry, we have the following inequalities:
\begin{enumerate}\itemsep=0pt
\item if $p \preceq q \preceq r$, the inverse triangle inequality holds:
\[
d(p,q) + d(q,r) \leq d(p,r),
\]
\item if $v$ and $w$ are timelike vectors, the inverse Cauchy--Schwartz inequality holds:
\[
\abs{g(v,w)} \geq \sqrt{- g(v,v)} \sqrt{- g(w,w)}.
\]
\end{enumerate}
\end{lemma}

\begin{proof}
Proofs can be found in \cite{Beem,ONeill}.
\end{proof}

Now we will try to adapt the Riemannian distance function by following the same procedure than above.

Let us consider a 4-dimensional oriented Lorentzian manifold $(M,g)$ and two points~$p$ and~$q$ on it such that $p \precc q$. Choose an arbitrary future directed timelike piecewise smooth curve $\gamma : [0,1] \rightarrow M$ with $\gamma(0) = p$ and $\gamma(1) = q$ (there must exist at least one such curve because $p \precc q$). Then, for each function $f \in C^\infty(M,\setR)$:
\begin{gather*}
f(q) - f(p) = f(\gamma(1)) - f(\gamma(0)) = \int_0^1 \dv{}{t} f(\gamma(t)) \, dt = \int_0^1 df(\dot\gamma(t)) \, dt = \int_0^1  g( \nabla f,\dot\gamma(t)) \, dt.
\end{gather*}
Notice that $\dot\gamma(t)$ is everywhere a future directed timelike vector. If we suppose that $\nabla f$ is everywhere timelike with constant orientation (we will take past directed), then $g( \nabla f,\dot\gamma(t))$ is of constant sign (hence positive):
\[
f(q) - f(p) = \int_0^1  g( \nabla f,\dot\gamma(t)) \, dt = \int_0^1 \abs{ g( \nabla f,\dot\gamma(t)) }\, dt.
\]


We can apply the inverse Cauchy--Schwartz inequality:
\begin{gather*}
\int_0^1  \abs{g( \nabla f,\dot\gamma(t))} \, dt  \geq \int_0^1  \sqrt{-g(\nabla f,\nabla f)} \sqrt{-g(\dot\gamma(t),\dot\gamma(t))} \, dt  \geq  \inf\set{{\sqrt{-g(\nabla f,\nabla f)}}}\; l(\gamma).
\end{gather*}

Now we can see that this construction holds also if $\gamma$ is causal instead of just timelike. Indeed, $g( \nabla f,\dot\gamma(t))$ is non-negative for a future directed lightlike vector $\dot\gamma(t)$ and Cauchy--Schwartz inequality becomes $g( \nabla f,\dot\gamma(t)) \geq 0$. So we can extend our construction to the case $p \prec q$.

Now we can see that $f(q) - f(p) \geq d(p,q)$ holds if we impose the following conditions on $f$:
\begin{itemize}\itemsep=0pt
\item  $\sup\, g( \nabla f, \nabla f ) \leq -1$;
\item $\nabla f$ is past directed.
\end{itemize}

The f\/irst condition is a global eikonal timelike condition. Once more, this condition is totally independent of any path on the manifold, so could give rise to a translation into an algebraic framework in order to be extended to noncommutative spaces. Of course this generalization is far from trivial and does not enter the scope of this paper.

The second condition is in reality not a restrictive condition but a simple choice of time orientation, which can be easily translated into a noncommutative framework. Indeed, functions obeying the f\/irst condition can be easily separated in two sets. To do that, just take two f\/ixed points $p_0$ and $q_0$ (or the corresponding states in a noncommutative algebra) such that $\inf_f \abs{ f(q_0) -f(p_0) } > 0$. Our two sets are
\begin{gather*}
\set{ f : g( \nabla f, \nabla f ) \leq -1,\; f(q_0) - f(p_0) > 0}\!\!\quad\text{and}\quad\!\!  \set{ f : g( \nabla f, \nabla f ) \leq -1,\; f(q_0) - f(p_0) < 0}\!
\end{gather*}
 and the choice of one of them corresponds to the choice of a particular time orientation.

Once more, we can extend the set of functions $f \in C^\infty(M,\setR)$ to a larger set, including continuous functions a.e.\ dif\/ferentiable, and use the weaker condition $\text{ess}\sup\, g( \nabla f, \nabla f ) \leq -1$. Nevertheless, due to the Lorentzian characteristic of the space, Lipschitz functions are not suitable. We must extend to an even larger set which we will denote here by $\mathcal A$. In fact, $\mathcal A$ must be a set of continuous functions on $M$ which are a.e.\ dif\/ferentiable and respect the fundamental theorem of calculus at least on causal paths (absolutely continuous functions).

There are many ways to characterize such set $\mathcal A$, and the choice of a particular one should take in account the generalization to noncommutative spaces. One can use the $\lambda$-absolutely continuous condition developed mainly by J.~Maly and S.~Hencl for functions of several variables or use Sobolov spaces $W^{1,p}(M)$ ($p>4$)~\cite{Hen,Mal}.

Another possibility for the space $\mathcal A$ is the space of continuous causal functions, i.e.\ functions which do not decrease along every causal future-directed curve. Indeed, by Lebesgue dif\/ferentiation theorem, such functions are of bounded variation on any future directed timelike smooth curve $\gamma : [0,1] \rightarrow M$, so a.e.\ dif\/ferentiable, and the fundamental theorem turns to be an inequality:
\[
f(q) - f(p) = f(\gamma(1)) - f(\gamma(0)) \geq \int_0^1 \dv{}{t} f(\gamma(t)) \, dt
\]
and the above construction remains valid.

From now, and for the rest of this paper, we will f\/ix $\mathcal A \subset C(M,\setR)$ as the space of continuous causal functions on $M$. Then we present our principal theorem, which takes place in a Lorentzian manifold with a globally hyperbolic condition:

\begin{theorem}\label{mainth}
Let $(M,g)$ be a globally hyperbolic spacetime and $d$ the Lorentzian distance function on $M$, then
\[
d(p,q) = \inf\set{ \langle f(q)-f(p) \rangle \ :\  f \in \mathcal A,\  \text{\rm ess} \sup g( \nabla f, \nabla f ) \leq -1, \ \nabla f \text{ past directed}}
\]
where $\langle \alpha \rangle = \max\set{0,\alpha}$.

\end{theorem}

In the next section, we give a proof of this theorem with $\mathcal A$ being the space of continuous causal functions, but the result is still valid with any suitable space $\mathcal A$, provided that the Lorentzian distance function belongs to $\mathcal A$ as a function of its second argument. The globally hyperbolic condition is a useful tool in order to construct functions giving the equality. The reason is that a global eikonal condition can be obtained only if we have a global knowledge on the causal structure of the manifold, and the globally hyperbolic condition does so.


\section{Proof of the Lorentzian distance function}


In the following, we will work in a globally hyperbolic spacetime $(M,g)$ (necessarily oriented). Our goal is to prove mathematically the Theorem \ref{mainth} above. We have already obtained the result:

\begin{lemma}\label{lemmainequa}
If $p \prec q$, then
\begin{gather*}
\forall \, f \in \mathcal A\text{ such that } \text{\rm ess} \sup g( \nabla f, \nabla f ) \leq -1   \text{ and }  \nabla f \text{ is past directed}, \  f(q)-f(p)   \geq  d(p,q).
\end{gather*}
\end{lemma}

We will separate the remaining in three cases:
\begin{itemize}\itemsep=0pt
\item if $p \prec q$, we will show that there exists a sequence of functions $f_\epsilon \in\mathcal A$ ($\epsilon > 0$) respecting the two given conditions and such that $d(p,q) \leq f_\epsilon(q) -f_\epsilon(p) < d(p,q) + \epsilon$ (Lemma \ref{equality});
\item if $p \succ q$, then there exists a function $f \in\mathcal A$ respecting the two given conditions and such that \mbox{$f(q)-f(p) \leq 0$} (Corollary \ref{equalityrev});
\item if $p \nprec q$ and $q \nprec p$, then there exists a sequence of functions $f_\epsilon \in\mathcal A$ ($\epsilon > 0$) respecting the two given conditions and such that $\abs{ f_\epsilon(q) -f_\epsilon(p) } < \epsilon$ (Lemma \ref{zero}).
\end{itemize}

With these three results, the proof of the Theorem \ref{mainth} will be completed, but we need to start f\/irst with some technical lemmas.

\begin{lemma}
In Lorentzian geometry, we have the following properties:
\begin{enumerate}\itemsep=0pt
\item if $(M,g)$ is globally hyperbolic, then $d$ is a finite continuous function on $M \times M$;

\item $d_p = d(p,\cdot) $ satisfies the timelike eikonal equation $g( \nabla d_p , \nabla d_p) = -1$ on $I^+(p) \setminus C^+(p)$ where $C^+(p)$ is called the cut locus of $p$;

\item if $(M,g)$ is globally hyperbolic, the cut locus  $C^+(p)$ has measure zero;

\item  if $(M,g)$ is globally hyperbolic, then $d_p\in\mathcal A$.
\end{enumerate}
\end{lemma}

\begin{proof}
Proof of (1) can be found in \cite{Beem,ONeill}, proof of (2) in~\cite{EGK} and (3) can be found in~\cite{Mor}. By def\/inition of the distance, $d_p$ is trivially non-decreasing along every causal future-directed curve.
\end{proof}

The basic idea of our construction is to create some functions with a dif\/ferent behaviour in two regions:
\begin{itemize}\itemsep=0pt
\item the f\/irst region is a region containing the points $p$ and $q$ where these functions correspond to a simple suitable distance function;
\item the second region is the remaining of the manifold where these functions are locally f\/inite sums of distance functions in order to have the eikonal condition respected in the whole space.
\end{itemize}

The next lemma will be useful to construct the second region.

\begin{lemma}\label{lemmaprincipal}
Let $S$ be a smooth spacelike Cauchy surface and two points $q \in J^+(S)$ $($resp.\ $q \in J^-(S))$ and $q' \in I^+(q)$ $($resp.\ $q' \in I^-(q))$. There exists a~function $f \in \mathcal A$ $($resp.\ $-f \in \mathcal A)$ such that:
\begin{itemize}\itemsep=0pt
\item $f \geq 0 $;
\item $g( \nabla f, \nabla f ) \leq -1$ and $\nabla f$ is past directed (resp.\ future directed) where $f > 0$, except on a~set of measure zero $($so $f$ is a.e.\ differentiable$)$;
\item $f > 0$ on $J^+(S) \setminus I^-(q')$ $($resp.\ on $J^-(S) \setminus I^+(q'))$;
\item $f = 0$ on $J^-(q)$ $($resp.\ on $J^+(q))$.
\end{itemize}
\end{lemma}

\begin{proof}
At f\/irst, notice that $J^-(q) \subset  I^-(q')$. Indeed, $q \precc q'$, so $I^-(q) \subset  I^-(q')$. To verify that $J^-(q) = \overline{I^-(q) }\subset  I^-(q')$ consider that $I^-(q') = \set{ p\in M : d(p,q') > 0}$ as shown in \cite{Beem} and that for every point $ z\in J^-(q)$, $d(z,q') \geq d(z,q) + d(q,q') > 0$.

On the smooth spacelike surface $S$ we can consider the Riemannian metric $g_R$ which is the restriction of the global metric $g$ to $S$, with a Riemannian distance $d_R$ and a topology associated. If we take a point $p \in I^-(S)$, the intersection $I^+(p) \cap S$ is an open subset of $S$ (in the topology of $S$) with a f\/inite diameter $d_R(I^+(p) \cap S) < \infty $  because $J^+(p) \cap J^-(S)$ is compact (see \cite{Wald}). We will work with points $p$ closed to the surface $S$ such that $d_R(I^+(p) \cap S)$ is small. Let us def\/ine
\[
P = \set{ p \in I^-(S) \setminus J^-(q) : d_R(I^+(p) \cap S) < 1}.
\]

Then the collection
\[
W = \set{ I^+(p) \cap S : p \in P}
\]
is an open covering of the closed set $S \setminus I^-(q')$ (because we have $I^-(q') \supset J^-(q)$). We will show that there exists a locally f\/inite subcovering by a method similar to \cite{BS03}. Let us f\/ix $s \in S$ and consider open and closed balls $B_s(r)$ and $\bar B_s(r)$ in~$S$ of center $s$ and radius $r$ for the distance~$d_R$. The following subsets are compact in $S$:
\[
S_n = \bar B_s(n) \setminus \parenthese{ B_s(n-1) \cup I^-(q')}, \quad n \in \setN \qquad \text{with} \qquad \bigcup_{n\in \setN} S_n = S \setminus I^-(q').
\]

For each $S_n$ we can f\/ind a f\/inite subset $\set{W_{1n}, \dots, W_{k_n n}} \subset W$ which covers $S_n$. Then $W' = \set{W_{kn} : n\in \setN , k = 1, \dots, k_n}$ is a locally f\/inite subcovering of $S \setminus I^-(q')$ because every $W_{kn}$ has diameter smaller that 1. So there exists a subset of points $P' \subset P$ such that $W' = \set{ I^+(p) \cap S : p \in P'}$ is locally f\/inite.

From that, we can show that $\set{I^+(p) : p \in P'}$ is a locally f\/inite covering of  $J^+(S) \setminus I^-(q')$. Indeed, take a point $z$ in $J^+(S) \setminus I^-(q')$. The set $I^-(z) \cap S$ is an open set of f\/inite diameter in $S$ which intersects only a f\/inite number of $W_{kn}$ (and must intersect at least one because $S \setminus I^-(q')$ is a Cauchy surface for $M \setminus I^-(q')$) and so $I^-(z)$ contains a non empty but f\/inite subset of $P'$. The same reasoning can be done for a small neighbourhood of $z$.

Now we can construct the non-negative function:
\[
f(z) = \sum_{p \in P'} d(p,z) = \sum_{p \in P'} d_p(z).
\]
This function is well def\/ined because the sum is pointwise f\/inite and is continuous by continuity of the distance function. $f$ is null on $J^-(q)$ because no $p \in P'$ belongs to $J^-(q)$ and is positive on $J^+(S) \setminus I^-(q')$ because every point in $J^+(S) \setminus I^-(q')$ is strictly inside the light cone of at least one $p \in P'$. For every $z \in M$ we can f\/ind a neighbourhood where $f$ is  a f\/inite sum $\sum d_p$ of distance functions, so $f\in\mathcal A$. Because $\nabla d_p$ is well def\/ined except on a set of measure zero, by countability of the measure, the locally f\/inite sum $\nabla f = \sum \nabla d_p$ is well def\/ined except on a~set of measure zero. Then where $\nabla f$ is well def\/ined and $f$ positive, we have that $\nabla f$ is timelike past directed (because it is the sum of null or timelike past directed vectors), and
\[
g( \nabla f, \nabla f ) =  \sum_{p} g( \nabla d_p, \nabla d_p ) +  2 \sum_{p \neq p'} g( \nabla d_p, \nabla d_{p'} ) \leq -1,
\]
where the f\/irst sum contains terms equal to $-1$ or $0$, with at least one term equal to $-1$, and the second sum contains terms negative or null because all $\nabla d_p$ are null or timelike past directed.

In the reverse case, ($q \in J^-(S)$ and $q' \in I^-(q)$), we can do an identical proof by reversing future and past sets and by taking $f(z) = \sum_{p \in P'} d(z,p)$ as a function with null or timelike future oriented gradient which is non-increasing along every causal future-directed curve.
\end{proof}

\begin{corollary}\label{corollaryprincipal}
Let $S$ be a smooth spacelike Cauchy surface and four points $q_1,q_2 \in J^+(S)$ and $q_1' \in I^+(q_1)$, $q_2' \in I^+(q_2)$ $($resp.\ $q_1,q_2 \in J^-(S)$, $q_1' \in I^-(q_1)$, $q_2' \in I^-(q_2))$. There exists a~function $f \in \mathcal A$ $($resp. $-f \in \mathcal A)$ such that:
\begin{itemize}\itemsep=0pt
\item $f \geq 0 $;
\item $g( \nabla f, \nabla f ) \leq -1$ and $\nabla f$ is past directed $($resp.\ future directed$)$ where $f > 0$, except on a~set of measure zero;
\item $f > 0$ on $J^+(S) \setminus \parenthese{I^-(q_1') \cup I^-(q_2')}$ $($resp.\ on $J^-(S) \setminus \parenthese{I^+(q_1') \cup I^+(q_2')})$;
\item $f = 0$ on $J^-(q_1) \cup J^-(q_2)$ $($resp.\ on $J^+(q_1) \cup J^+(q_2))$.
\end{itemize}
\end{corollary}

\begin{proof}
The proof is identical to Lemma \ref{lemmaprincipal} except that we start with the set \[
P = \set{ p \in I^-(S) \setminus \parenthese{J^-(q_1) \cup J^-(q_2)} : d_R(I^+(p) \cap S) < 1}
\]
to create a locally f\/inite covering of $J^+(S) \setminus \parenthese{ I^-(q_1') \cup I^-(q_2') }$.
\end{proof}

The remaining of this section consists of the construction of the dif\/ferent functions.

\begin{lemma}\label{equality}
If $p \prec q$, then $\forall\, \epsilon > 0$ there exists a function $f\in \mathcal A$ such that:
\begin{itemize}
\item $\text{\rm ess} \sup g( \nabla f, \nabla f ) \leq -1$;
\item $\nabla f$ is past directed;
\item $f(q)-f(p) \geq 0$;
\item $\abs{(f(q)-f(p)) - d(p,q)} < \epsilon$.
\end{itemize}

\end{lemma}

\begin{proof}
Let us choose a smooth spacelike Cauchy surface $S$ containing $q$. Such smooth spacelike surface exists as shown by works of Bernal and S\'anchez \cite{BS03,BS04}. Then let us choose two free points~$p'$ and~$q'$ such that $p' \in I^-(p)$ and $q' \in I^+(q)$.  Because $q \in I^+(p')$ and $q\in S$ we can choose $q'$ close to $q$ such that $I^-(q') \cap J^+(S) \subset I^+(p')$.

We can apply the Lemma \ref{lemmaprincipal} to $S$, $q$ and $q'$ to get a function $f_1$ with properties:
\begin{itemize}\itemsep=0pt
\item $f_1 \in \mathcal A$;
\item $f_1 \geq 0 $;
\item $g( \nabla f_1, \nabla f_1 ) \leq -1$ and $\nabla f_1$ is past directed where $f_1 > 0$, except on a set of measure zero;
\item $f _1> 0$ on $J^+(S) \setminus I^-(q')$;
\item $f_1 = 0$ on $J^-(q)$,
\end{itemize}
and then to $S$, $p$ and $p'$ to get a function $f_2$ with properties:
\begin{itemize}\itemsep=0pt
\item $-f_2 \in \mathcal A$;
\item $f_2 \geq 0 $;
\item $g( \nabla f_2, \nabla f_2 ) \leq -1$ and $\nabla f_2$ is future directed where $f_2> 0$, except on a set of measure zero;
\item $f_2 > 0$ on $J^-(S) \setminus I^+(p')$;
\item $f_2 = 0$ on $J^+(p)$.
\end{itemize}

Then the function:
\[
f_0 = f_1 - f_2
\]
has the following properties:
\begin{itemize}\itemsep=0pt
\item $f_0\in\mathcal A$;
\item $f_0 = 0$ on the compact $J^-(q) \cap J^+(p)$;
\item the support of $\nabla f_0$ includes the set $\parenthese{J^+(S) \setminus I^-(q')} \cup  \parenthese{J^-(S) \setminus I^+(p')}$;
\item $g( \nabla f_0, \nabla f_0) \leq -1$ and $\nabla f_0$ is past directed on its support, except on a set of measure zero.
\end{itemize}
The last assertion comes from
\[
g( \nabla f_0, \nabla f_0 ) =  g( \nabla f_1, \nabla f_1 ) + g( \nabla f_2, \nabla f_2 ) -  2 g( \nabla f_1, \nabla f_2 ) \leq -1,
\]
where the f\/irst term is equal to $-1$ on the support of $f_1$ and is non-positive elsewhere, the second term is equal to~$-1$ on the support of $f_2$ and is non-positive elsewhere, and the last term is non-positive because $\nabla f_1$ and $\nabla f_2$ are not in the same orientation.

We can now def\/ine the function:
\[
f = f_0 + d_{p'}  \quad \Lequi \quad f(z) = f_0 + d(p',z)
\]
which has the following properties:
\begin{itemize}\itemsep=0pt
\item $f\in\mathcal A$;
\item $f(p)= d(p',p)$;
\item $f(q)= d(p',q)$;
\item the support of $\nabla f$ is $M$;
\item $g( \nabla f, \nabla f) \leq -1$ and $\nabla f$ is past directed on $M$, except on a set of measure zero.
\end{itemize}

To verify that the support of $\nabla f$ is $M$, we can see that $M \setminus J^+(p') \subset \parenthese{J^+(S) \setminus I^-(q')} \cup  \parenthese{J^-(S) \setminus I^+(p')} $ (recall that $I^-(q') \cap J^+(S) \subset I^+(p')$) and that the support of $d_{p'}$ is $J^+(p')$.

So now we have a function $f$ such that $f(q)-f(p) = d(p',q) - d(p',p) \geq d(p,q) \geq 0$ by the inverse triangle inequality.

Let us set the function $\alpha(p') = \parenthese{f(q) - f(p)} - d(p,q) = \parenthese{d(p',q) - d(p',p)} - d(p,q)$. $\alpha$~is a~continuous function because the distance function is continuous, and $\alpha(p) = 0$. So it is always possible to choose the initial point $p' \in I^-(p)$ such that $\abs{\alpha(p')} < \epsilon$.
\end{proof}

\begin{corollary}\label{equalityrev}
If $p \succ q$, then there exists a function $f\in \mathcal A$ such that:
\begin{itemize}\itemsep=0pt
\item $\text{\rm ess} \sup g( \nabla f, \nabla f ) \leq -1$;
\item $\nabla f$ is past directed;
\item $f(q)-f(p) \leq 0$.
\end{itemize}
In particular, $\langle f(q)-f(p) \rangle = 0$.

\end{corollary}

\begin{proof}
This is trivial by switching $p$ and $q$ in Lemma \ref{equality}.
\end{proof}

\begin{lemma}\label{zero}
If $p \nprec q$ and $q \nprec p$, then $\forall \, \epsilon > 0$ there exists a function $f\in \mathcal A$ such that:
\begin{itemize}
\item $\text{\rm ess} \sup g( \nabla f, \nabla f ) \leq -1$;
\item $\nabla f$ is past directed;
\item $\abs{f(q)-f(p)} < \epsilon$.
\end{itemize}
\end{lemma}

\begin{proof}
The case $p=q$ is trivial, so we will suppose $q\notin J^{\pm}(p)$.

Let us choose a smooth spacelike Cauchy surface $S$ containing $q$ and assume $p \in J^-(S)$ (otherwise we exchange the role of $p$ and $q$). Then we choose the following free points:
\begin{itemize}\itemsep=0pt
\item $\tilde p \in S$ such that $\tilde p \in J^+(p)$ (note: $\tilde p = p$ if $p\in S$);
\item $p' \in I^-(p)$ and $q' \in I^-(q)$ such that the sets $J^+(p') \cap S$ and $J^+(q') \cap S$ are disjoint (this is always possible because $q\in S$ and $q\notin J^{+}(p)$);
\item $\tilde p' \in I^+(\tilde p)$ and $\tilde q' \in I^+(q)$ such that $I^-(\tilde p') \cap J^+(S) \subset I^+(p')$ and $I^-(\tilde q') \cap J^+(S) \subset I^+(q')$.
\end{itemize}
Then we have an analogous situation than in Lemma \ref{equality} but with two disjoint sets
\[
\parenthese{J^+(S) \cap I^-(\tilde p')} \cup  \parenthese{J^-(S) \cap I^+(p')}\qquad\text{and}\qquad\parenthese{J^+(S) \cap I^-(\tilde q')} \cup  \parenthese{J^-(S) \cap I^+(q')}
\]
the f\/irst one containing the compact $J^-(\tilde p) \cap J^+(p)$ and the second the point $q$.

We can then apply the Corollary \ref{corollaryprincipal} a f\/irst time to $S$, $\tilde p$, $q$, $\tilde p'$ and $\tilde q'$ to get a function $f_1$ with null of past directed gradient, and a second time to $S$, $p$, $q$, $p'$ and $q'$ to get a function $f_2$ with null of future directed gradient. In the same way as in Lemma \ref{equality} we f\/ind a function:
\[
f_0 = f_1 - f_2
\]
with the following properties:
\begin{itemize}\itemsep=0pt
\item $f_0\in\mathcal A$;
\item $f_0 = 0$ on the compact $J^-(\tilde p) \cap J^+(p)$ and on $q$;
\item the support of $\nabla f_0$ includes the set $\parenthese{J^+(S) \!\setminus \! \parenthese{I^-(\tilde p') \cup I^-(\tilde q')}} \cup    \parenthese{J^-(S)\! \setminus \!\parenthese{I^+(p') \cup I^+(q') }}$;
\item $g( \nabla f_0, \nabla f_0) \leq -1$ and $\nabla f_0$ is past directed on its support, except on a set of measure zero.
\end{itemize}

Finally we def\/ine the function $f$:
\[
f = f_0 + d_{p'} + d_{q'}  \quad \Lequi \quad   f(z) = f_0 + d(p',z) + d(q',z)
\]
which has the following properties:
\begin{itemize}\itemsep=0pt
\item $f\in\mathcal A$;
\item $f(p)= d(p',p)$;
\item $f(q)= d(q',q)$;
\item the support of $\nabla f$ is $M$;
\item $g( \nabla f, \nabla f) \leq -1$ and $\nabla f$ is past directed on $M$, except on a set of measure zero.
\end{itemize}

Once more we can choose $p' \in I^-(p)$ and $q' \in I^-(q)$ such that $d(p',p) < \frac\epsilon2$, $d(q',q) < \frac\epsilon2$ and conclude that $\abs{f(q)-f(p)} < \epsilon$.
\end{proof}

Lemma \ref{zero} concludes the proof of the Theorem~\ref{mainth}.

Notes:
\begin{itemize}\itemsep=0pt
\item For every set $\mathcal A \subset C(M,\setR)$ of a.e. dif\/ferentiable functions with absolute continuity condition respected on causal paths and such that $d_p \in \mathcal A$, the construction above holds. Indeed, path integrations take place into compact sets and the f\/inal function $f$ constructed in Lemmas~\ref{equality} and~\ref{zero} is a f\/inite sum of distance functions on each compact set, so the fundamental theorem of calculus still holds. A proof of the absolute continuity of the function $d_p$ would be then a good improve of our result.
\item We think that the f\/inal function $f$ in Lemmas \ref{equality} and \ref{zero} could also be replaced by the distance $f(z) = d(S,z) - d(z,S)$ to a suitable smooth Cauchy surface $S$ (in fact, a surface containing~$p$ and such that the geodesic from~$p$ to~$q$ is maximal in Lemma \ref{equality} or con\-taining both~$p$ and~$q$ in Lemma \ref{zero}) if one can prove the existence of such surface, the a.e.\ dif\/ferentiability and the absolute continuity of the distance $d(S,z)$.
\end{itemize}



\begin{thebibliography}{99}
\bibitem[1]{Beem}
Beem J.K., Eherlich P.E., Easley K.L.,
Global Lorentzian geometry,  2nd ed., {\it Monographs and Textbooks in Pure and Applied Mathematics}, Vol.~202, Marcel Dekker, Inc., New York, 1996.

\bibitem[2]{BS03}
Bernal A.N., S\'anchez M.,
On smooth Cauchy hypersurfaces and Geroch's splitting theorem,
\href{http://dx.doi.org/10.1007/s00220-003-0982-6}{{\it Comm. Math. Phys.}} {\bf 243} (2003), 461--470,
\href{http://www.arxiv.org/abs/gr-qc/0306108}{gr-qc/0306108}.

\bibitem[3]{BS04}
Bernal A.N., S\'anchez M.,
Smoothness of time functions and the metric splitting of globally hyperbolic spacetimes,
\href{http://dx.doi.org/10.1007/s00220-005-1346-1}{{\it Comm. Math. Phys.}} {\bf 257} (2005), 43--50,
\href{http://www.arxiv.org/abs/gr-qc/0401112}{gr-qc/0401112}.

\bibitem[7]{EGK}
Erkeko\u{g}lu F., Garc\'ia-R\'io E., Kupeli D.N.,
On level sets of Lorentzian distance function,
\href{http://dx.doi.org/10.1023/A:1025779017980}{{\it Gen. Relativity Gravitation}} {\bf 35} (2003), 1597--1615.

\bibitem[9]{Hen}
Hencl S.,
On the notions of absolute continuity for functions of several variables,
\href{http://dx.doi.org/10.4064/fm173-2-5}{{\it Fund. Math.}} {\bf 173} 2002, 175--189.

\bibitem[10]{Mal}
Mal\'y J.,
Absolutely continuous functions of several variables,
\href{http://dx.doi.org/10.1006/jmaa.1998.6246}{{\it J.~Math. Anal. Appl.}} {\bf 231} (1999), 492--508.

\bibitem[11]{Mor}
Moretti V.,
Aspects of noncommutative Lorentzian geometry for globally hyperbolic spacetimes,
\href{http://dx.doi.org/10.1142/S0129055X03001886}{{\it Rev. Math. Phys.}} {\bf 15} (2003), 1171--1217,
\href{http://arxiv.org/abs/gr-qc/0203095}{gr-qc/0203095}.

\bibitem[12]{ONeill}
O'Neill B.,
Semi-Riemannian geometry. With applications to relativity,
 {\it Pure and Applied Mathematics}, Vol.~103, Academic Press, Inc., New York, 1983.

\bibitem[15]{Wald}
Wald R.M.,
General relativity,  University of Chicago Press, Chicago, IL, 1984.

\end{thebibliography}
\end{document}
